\section{A finite-overlap physical atlas and original-roof source height}
\label{sec:at-atlas}
We prove a height estimate for the complete class having a fixed
margin in the \emph{original} incidence and clearance tests.  The
class includes both critical and noncritical physical trajectories.
It is not restricted to trajectories that continue to a critical
orbit.  The argument uses actual two-dimensional families of
trajectories and bounds their overlap before summing their positive
optical weights.  Neither a whole-cylinder pressure estimate nor a
mixing theorem is used.

Write $r_g=q_0^{1/4}$, where $q_0=47/53$ remains unchanged.
For $0<\varepsilon\le1/4$ let $\mathcal G_{\varepsilon,m,R}$ be
the regular physical records satisfying
\begin{equation}\label{eq:at-source}
 |p_0|,|p_m|\le5/6,\qquad
 c_i\ge\varepsilon r_g^{d_i}\ (0\le i\le m),\qquad
 \Delta_i\ge\varepsilon r_g^{d'_i}\ (0\le i<m).
\end{equation}
These are source restrictions used in an estimate; the original
probability, event and microscopic windows are unchanged.  The
endpoint cap includes the entire original exact event.  Indeed,
$q_R<53/100$ and $q\mapsto2q/(1+q^2)$ is increasing on $(0,1)$,
so the largest section momentum is less than
$10600/12809+1/200<5/6$.  The other three types of section rectangle
have smaller momentum bounds.  No section-distance margin is
imposed in \eqref{eq:at-source}.

Keep the already proved constants $\delta_*,C_{\rm inc}$ from
\eqref{eq:pt-tube-constant}, $D_*,C_{\rm dist}$ from
\eqref{eq:pt-dist-constants}, $J_*,C_c,C_\Delta$ from
Section~\ref{sec:el-localization}, and put
\begin{equation}\label{eq:at-constants}
 g=9/40,\qquad U=320/9,\qquad
 a_\varepsilon=\delta_*\varepsilon/16,\qquad
 \rho_\varepsilon^{\rm at}=\frac{g a_\varepsilon}{8R_+},\qquad
 C_{\rm at}=\frac{25U C_{\rm dist}^2}
                   {2c_*(\rho_\varepsilon^{\rm at})^2}.
\end{equation}
We abbreviate $\rho=\rho_\varepsilon^{\rm at}$ within this section.
It is an endpoint-momentum chart radius, never the physical roof or
the microscopic comparison width $r_B$.  All constants and chart
radii are fixed before the collision count tends to infinity.

\begin{lemma}[A common momentum neighborhood of every guarded record]
\label{lem:at-ball}
For $x\in\mathcal G_{\varepsilon,m,R}$ put
$q_x=(\alpha_0,\alpha_m)$ in local endpoint lifts and
$\zeta_x=(p_0,p_m)$.  The actual endpoint ball
$|q-q_x|_2<a_\varepsilon$ is physical with the same lifted word,
and all its incidences and clearances satisfy the half-margins
\eqref{eq:pt-half-margin}.  Its endpoint incidences exceed $1/2$.
Every $|\zeta-\zeta_x|<4\rho$ has a unique actual inverse on
that ball, with
\begin{equation}\label{eq:at-inverse-ball}
 |q(\zeta)-q_x|_2\le R_+|\zeta-\zeta_x|/g<a_\varepsilon/2.
\end{equation}
Here and below uniqueness also holds among all regular physical
trajectories of the same lifted word with the given $\zeta$.
\end{lemma}
\begin{proof}
On the closed ball, $|h_0|+|h_m|\le2a_\varepsilon
=\delta_*\varepsilon/8<\delta_*\varepsilon$.
Theorem~\ref{thm:pt-physical} gives the actual trajectories and all
half-margins, on a neighborhood of this closed ball.  Its signed
tangential estimate bounds each endpoint momentum change by
$C_{\rm inc}\delta_*\varepsilon/8\le\varepsilon/16\le1/64$.
Thus $|p_0|,|p_m|\le5/6+1/64=163/192<17/20$ and
$c_0,c_m\ge\sqrt{111}/20>1/2$.
The original second variation therefore gives $G=D_q^2F\succeq gI$.

For the desired momentum pair minimize
$F(q)-RD\zeta\cdot(q-q_x)$ on the closed ball, where
$D=\operatorname{diag}(-1,1)$.  At its boundary, putting
$h=q-q_x$ and using $\nabla F(q_x)=RD\zeta_x$, the outward
radial derivative is at least
\[
 g a_\varepsilon^2-R_+|\zeta-\zeta_x|a_\varepsilon
 >g a_\varepsilon^2/2>0.
\]
A boundary minimizer is impossible.  Strong convexity gives one
interior minimizer, where $\nabla F=RD\zeta$.
The inverse function theorem and strong monotonicity give its
smooth dependence and \eqref{eq:at-inverse-ball}.  This minimization
takes place inside a ball already proved physical; it does not
assert that an arbitrary stationary continuation is physical.
Finally, with equal endpoint momenta the left side of the support
inequality \eqref{eq:v111-support-sum} is zero.  Strict positivity
of every physical incidence forces all contacts to coincide.
This proves the asserted global wordwise uniqueness, even across
disconnected original event components.
\end{proof}

\begin{lemma}[Physical tubes with multiplicity at most twenty-five]
\label{lem:at-cover}
For each fixed $R,m,\varepsilon$ there is a finite family
$\{T_{w,z}\}$ of actual physical source sets, indexed by lifted
words $w$ and points $z$ of the grid $\rho\mathbb Z^2$, with
these properties.
Every $T_{w,z}$ is parametrized by the \emph{whole} momentum disk
$B(z,2\rho)$, carries word $w$, and contains a unique physical
center $x_{w,z}$ with momentum $z$.  Their union covers
$\mathcal G_{\varepsilon,m,R}$.  Each physical initial point is
in at most $25$ of the tubes, counting all words and grid points.
Writing
$u(x)=|\partial_{\alpha_0}p_m(x)|^{-1}$, one has throughout a tube
\begin{equation}\label{eq:at-weight}
 C_{\rm dist}^{-1}u(x_{w,z})\le u(x)\le
 C_{\rm dist}u(x_{w,z}),\qquad
 \nu(T_{w,z})\ge\frac{\rho^2}{C_{\rm dist}}u(x_{w,z}).
\end{equation}
For any two points $x,y$ in a tube,
\begin{align}
 |\Delta_i(x)-\Delta_i(y)|
     &\le2C_\Delta a_\varepsilon\sigma^{d'_i},
                                   &&0\le i<m,\label{eq:at-clearance}\\
 |c_i(x)-c_i(y)|
     &\le2C_c a_\varepsilon\sigma^{d_i},
                                   &&0\le i\le m,\label{eq:at-incidence}\\
 d_\infty(T_R^ix,T_R^iy)
     &\le2C_c a_\varepsilon\sigma^{d_i},
                                   &&0\le i\le m.\label{eq:at-state}
\end{align}
The states in the last line are actual collision states, including
the incoming definition of the terminal tangential momentum.
\end{lemma}
\begin{proof}
Assign each $\zeta_x$ of a source point to a nearest grid point,
with a fixed convention on ties.  The Euclidean distance is at
most $\rho/\sqrt2<\rho$.  For every nonempty pair $(w,z)$ choose
one such source point $x$.  There are finitely many relevant grid
points and finitely many physical words at fixed $m$, by the
proved physical horizon.  Apply Lemma~\ref{lem:at-ball} at $x$.
Since $B(z,2\rho)\subset B(\zeta_x,3\rho)$, its inverse gives
the specified tube and center.  The local parametrization is
injective as a physical source map: its momenta determine the
point by wordwise uniqueness.  Every other assigned source point
with the same pair $(w,z)$ agrees with the constructed inverse at
its own momentum by that same global uniqueness.  Thus the family
covers the entire class, not just the chosen representatives.
Different constructions for a pair define exactly the same
physical inverse on the disk.

A physical initial point determines its unique complete lifted
word.  If it lies in $T_{w,z}$, its momentum is within $2\rho$ of
$z$.  Each coordinate interval of length $4\rho$ contains at
most five points of $\rho\mathbb Z$.  There are consequently at
most $25$ possible grid centers.  This multiplicity bound uses the
physical momentum map; it introduces neither independent marks
nor a factor equal to the number of collision words.

In the chosen endpoint ball all tube points and the center have
distance less than $a_\varepsilon/2$ from $q_x$.  Any two can be
joined by their straight segment inside this actual physical ball,
of endpoint $\ell^\infty$ length at most $a_\varepsilon$.
All along it $c_i\ge(\varepsilon/2)r_g^{d_i}$ and the endpoint
incidences exceed $1/2$.  The absolute logarithmic derivative
bound in Lemma~\ref{lem:pt-distortion}, integrated in both
directions, gives distortion at most
$\exp(D_*a_\varepsilon/\varepsilon)
\le\exp(D_*\delta_*/2)=C_{\rm dist}$.
In momentum coordinates the \emph{unconditioned} invariant
collision probability is $u\,d\zeta/(4\pi)$.  The disk area is
$4\pi\rho^2$.  Integrating the lower distortion bound proves
the volume estimate in \eqref{eq:at-weight}.

For clearance and incidence integrate
\eqref{eq:el-clearance} and \eqref{eq:el-incidence} along the
same physical segment, using $s_i\le2\sigma^{d_i}$.
The tangential derivative obeys the identical incidence bound:
differentiate $p_i=v_i\cdot t_i$, or
$p_m=v_{m-1}\cdot t_m$ at the terminal contact, and use the
direction and contact estimates of
Proposition~\ref{prop:el-observables}.  Also
$|D\alpha_i|\le J_*s_i\|\eta\|_\infty$ and $J_*\le C_c$.
This proves \eqref{eq:at-state} in the circular coordinate metric.
The clearance estimate is for the complete closed flight segment
and every third disk; nearest-disk switches and disks visited at
other times are already included in its Lipschitz proof.
\end{proof}

\begin{lemma}[The roof Jacobian on the entire endpoint cap]
\label{lem:at-roof-jacobian}
At every regular point in these tubes, write
$\zeta=s e_\theta\ne0$ and
$\widehat G=R^{-2}G$.  Then
\begin{equation}\label{eq:at-roof-jacobian}
 \frac1U\le h(x,\theta):=
 e_\theta^{\mathsf T}D\widehat G^{-1}De_\theta
 \le 2R_+,\qquad
 \frac{dL(x(se_\theta))}{ds}=s h(x,\theta),\qquad
 b(x)=\frac{u(x)}{4\pi c_*h(x,\theta)}.
\end{equation}
The last expression is the original section-conditioned source
density in $(t,\theta)$, with section and all exact-event tests
retained as indicators.  A momentum disk need not be centered at
the origin for this formula.
\end{lemma}
\begin{proof}
The Schur-complement proof of Lemma~\ref{lem:rc-metric} gives
$\widehat G\preceq(320/9)I$ for all $|p_0|,|p_m|<1$.
For $m\ge2$ its better bound is $170/9$; the larger constant
also treats the single-flight mixed entry.  Its lower bound uses
only $c_0,c_m\ge1/2$, so
$\widehat G\succeq(2R_+)^{-1}I$.  These componentwise caps
are enough; we do not require $|\zeta|_2<1$.
The exact identities
$\nabla_qF=RD\zeta$ and $D_\zeta q=RG^{-1}D$ now give the
radial derivative.  The inverse source Jacobian is $u$, polar
area is $s\,ds\,d\theta$, and $dt=s h\,ds$.  Cancellation
gives the displayed density, with precisely one $c_*^{-1}$ and
no omitted radius factor.  On each open radial interval in a
tube it is ordinary one-dimensional substitution.  For a fixed
word, Theorem~\ref{thm:v111-unique} prevents two distinct regular
points with the same roof and angle, even if an exact-event cut
disconnects that interval.  At $\zeta=0$ there is at most one
regular physical point per word; these finitely many source
points have zero collision area.  No singular endpoint trace is
discarded by this area statement.
\end{proof}

For a Borel physical source set $E$, let
$\operatorname{Sat}_\varepsilon(E)$ be the union of precisely
those tubes above which meet $E\cap\mathcal G_{\varepsilon,m,R}$.
This is a finite union of open physical source sets and is Borel
at each fixed radius and collision count.  The choices need not
be measurable jointly in $R$: the required order is
$\sup_R\esssup_t$, with radius-dependent null sets.  This
auxiliary union is used only on the upper side of an inequality;
it is never a conditioning event.

\begin{theorem}[Complete physical source height by tube-volume payment]
\label{thm:at-height}
For every $m\ge1$, $R\in[R_-,R_+]$, original exact labels $n,k$,
Borel $E$ and $0\le h\le1$, the original full-roof measure
\[
 \mu_E=(L_{m,R})_*
  (\one_E\one_{\mathcal G_{\varepsilon,m,R}}
                         \one_{\Pi_{n,k,m,R}}h\,d\nu_R^*)
\]
is absolutely continuous and satisfies
\begin{equation}\label{eq:at-height}
 \left\|\frac{d\mu_E}{dt}\right\|_\infty
 \le C_{\rm at}\,
       \nu\bigl(\operatorname{Sat}_\varepsilon(E)\bigr).
\end{equation}
In particular the full guarded physical source has height at most
$C_{\rm at}=O(\varepsilon^{-2})$, uniformly in all collision
counts, radii, labels and bounded selectors.  This conclusion has
no premise from Input~\ref{input:v112-geometric-pressure}, and
does not use the endpoint-kernel constant $C_{\rm mom}$.
\end{theorem}
\begin{proof}
Assign every source point in the covered class to the first tube
containing it in a fixed enumeration.  This is a Borel partition
of the original class.  Retain the assignment, $E$, the complete
event and $h$ in the coarea integrand.  Lemma~\ref{lem:at-roof-jacobian}
and the upper bound in \eqref{eq:at-weight} dominate that
integrand on its assigned tube by
$U C_{\rm dist}u(x_{w,z})/(4\pi c_*)$.
Only tubes meeting $E\cap\mathcal G_{\varepsilon,m,R}$ can
contribute.  For every such tube its lower volume bound gives
$u(x_{w,z})\le C_{\rm dist}\nu(T_{w,z})/\rho^2$.
The proved multiplicity bound, applied to this subfamily, yields
\[
 \sum_{(w,z):\,T_{w,z}\cap E\cap\mathcal G_\varepsilon\ne\varnothing}
 u(x_{w,z})
 \le\frac{25C_{\rm dist}}{\rho^2}
             \nu(\operatorname{Sat}_\varepsilon(E)).
\]
Integrating the uniform coarea bound over an angle of length
$2\pi$ gives exactly \eqref{eq:at-height}.

To justify absolute continuity before claiming its bound, apply
local substitution to nonnegative roof tests on relatively
compact regular charts, disjointize only their representations,
and exhaust them monotonically as in
Proposition~\ref{prop:v112-exhaustion}.  The estimates just
proved bound the completed positive sum.  All physical words
and all original cuts are included before taking its essential
supremum.  The underlying tube volume used $\nu$ without an
initial-section restriction; the source coarea used $\nu_R^*$
with the original event.  This explains the single factor
$c_*^{-1}$ and involves no change of probability.
\end{proof}

\begin{corollary}[Observable neighborhoods in the actual collision process]
\label{cor:at-observables}
Let $E$ be the union, over any designated contact or flight
indices, of the actual conditions
$\Delta_i\le\eta_i$, $c_i\le\xi_i$, and
$\operatorname{dist}_\infty(T_R^ix,\partial Y_R^*)\le\zeta_i$.
All widths are nonnegative.  Replace them respectively by
\begin{equation}\label{eq:at-widening}
 \eta_i^+=\eta_i+2C_\Delta a_\varepsilon\sigma^{d'_i},\quad
 \xi_i^+=\xi_i+2C_c a_\varepsilon\sigma^{d_i},\quad
 \zeta_i^+=\zeta_i+2C_c a_\varepsilon\sigma^{d_i},
\end{equation}
and denote the resulting original-observable union by $E^+$.
Then $\operatorname{Sat}_\varepsilon(E)\subset E^+$.
When $\eta_i^+<3/200$, $\xi_i^+\le1$, and
$\zeta_i^+\le1/200$, one has the complete bound
\begin{equation}\label{eq:at-observable-height}
 \left\|\frac{d\mu_E}{dt}\right\|_\infty
 \le C_{\rm at}
 \left\{\frac1{R_-}\sum_i\eta_i^+
              +\sum_i(\xi_i^+)^2
              +\frac1{2\pi}\sum_i\zeta_i^+\right\}.
\end{equation}
Each sum is only over its designated actual indices; overlap of
different witnesses requires no independence assumption.
\end{corollary}
\begin{proof}
A tube meeting $E\cap\mathcal G_\varepsilon$ has a point with
at least one of the stated witnesses.  Every point of that whole
tube has the corresponding widened witness, by
\eqref{eq:at-clearance}--\eqref{eq:at-state} and the
one-Lipschitz property of distance to the fixed original section
boundary.  This proves the inclusion without selecting the same
witness on different tubes.  The exact physical clearance flux
bounds $\nu\{\Delta_i\le u\}\le u/R_-$ in the stated range.
The invariant incidence marginal is
$\nu\{c_i\le u\}=1-\sqrt{1-u^2}\le u^2$.
Lemma~\ref{lem:ec-collar} bounds the section collar by $u/(2\pi)$.
Apply invariance of the actual collision map and the union bound,
then Theorem~\ref{thm:at-height}.  All tests remain on actual
physical flights and contacts.  The widened sets occur only in
this proved upper estimate; the original event and its
microscopic widths are not widened.
\end{proof}

The finite-overlap mechanism is a quantitative comparison between
a point's optical weight and a positive volume of actual nearby
trajectories.  It therefore supplies information absent from a
total-mass bound alone.  It does not show that the complement of
$\mathcal G_{\varepsilon,m,R}$ has small normalized roof height.
The factor $\varepsilon^{-2}$ is explicit, so no uniform passage
to zero guard margin is implicit in this theorem.
